\section*{الفصل \ref{ch:b2:huckel} --- نظرية هوكل والأنظمة المترافقة}

\begin{solution}{exo:b2:huckel:1}
مستويات الأليل ($n = 3$): $\alpha + 1.414\beta$، $\alpha$، $\alpha - 1.414\beta$. الكاتيون (إلكترونان):
$2\alpha + 2.828\beta$؛ والجذر (3): $3\alpha + 2.828\beta$؛ والأنيون (4): $4\alpha + 2.828\beta$. فالإلكترونات
الزائدة تذهب إلى المستوى غير الرابط ولا تغيّر إلا الجزء $\alpha$.
\end{solution}

\begin{solution}{exo:b2:huckel:2}
$2(\alpha + 1.618\beta) + 2(\alpha + 0.618\beta) = 4\alpha + 4.472\beta$.
\end{solution}

\begin{solution}{exo:b2:huckel:3}
البنزين (6)، وأنيون البنتاديينيل الحلقي (6)، والتروبيليوم (6): \omterm{def:b2:huckel:aromatic}{عطرية}. وكاتيون البنتاديينيل
الحلقي (4)، والبوتاديين الحلقي (4)، والأوكتاتتراين الحلقي المستوي (8): ضد \omterm{def:b2:huckel:aromatic}{عطرية} إذا كانت مستوية
(والأوكتاتتراين الحلقي الحقيقي على شكل حوض وغير \omterm{def:b2:huckel:aromatic}{عطري}).
\end{solution}

\begin{solution}{exo:b2:huckel:4}
مثمّن داخل الدائرة: المستويات $\alpha + 2\beta$، $\alpha + 1.414\beta$ (مرتين)، $\alpha$ (مرتين)،
$\alpha - 1.414\beta$ (مرتين)، $\alpha - 2\beta$. ثمانية إلكترونات: اثنان في الأدنى، وأربعة في
الزوج الأول، وواحد في كل مدار من الزوج غير الرابط.
\end{solution}

\begin{solution}{exo:b2:huckel:5}
$4.472\beta - 4\beta = 0.472\beta$، أي $0.472 \times 75 = \qty{35}{kJ/mol}$، مقابل نحو
\qty{10}{kJ/mol} من معطيات الهدرجة: فنموذج هوكل، مع $\beta$ موفَّقة على البنزين،
يبالغ في تقدير ترافق السلاسل المفتوحة.
\end{solution}

\begin{solution}{exo:b2:huckel:6}
$p_{12} = p_{34} = 0.894$، $p_{23} = 0.447$: الرابطتان الطرفيتان C1--C2 وC3--C4 هما
الأقصر، والرابطة المركزية أطول لكنها أقصر من رابطة أحادية.
\end{solution}

\begin{solution}{exo:b2:huckel:7}
الأنيون، ستة إلكترونات: $2 \times 2 + 4 \times 0.618$، $E_\pi = 6\alpha + 6.472\beta$، طبقة مغلقة:
\omterm{def:b2:huckel:aromatic}{عطري}. والكاتيون، أربعة: $2 \times 2 + 2 \times 0.618$، $4\alpha + 5.236\beta$، بإلكترون واحد
في كل مدار من الزوج المتساوي الطاقة: \omterm{def:b2:huckel:aromatic}{ضد عطري}، وغير مستقر جدًا.
\end{solution}

\begin{solution}{exo:b2:huckel:8}
الحلقة السباعية: $\alpha + 2\beta$، $\alpha + 1.247\beta$ (مرتين)، ثم مستويات مضادة للربط. ستة
إلكترونات، كما في الكاتيون \ce{C7H7+}، تملأ المستويات الرابطة تمامًا: كاتيون
\omterm{def:b2:huckel:aromatic}{عطري}. فيتأين المركب بسهولة إلى تروبيليوم وبروميد.
\end{solution}

\begin{solution}{exo:b2:huckel:9}
$1.802 + 1.247 + 0.445 - 0.445 - 1.247 - 1.802 = 0$: مجموع الطاقات هو $6\alpha$.
\end{solution}

\begin{solution}{exo:b2:huckel:10}
$n = 6$: $m_k = 2\cos(k\pi/7) = 1.802$، $1.247$، $0.445$، $-0.445$، $-1.247$، $-1.802$. ستة
إلكترونات: $E_\pi = 6\alpha + 2(1.802 + 1.247 + 0.445)\beta = 6\alpha + 6.988\beta$؛
\omterm{def:b2:huckel:delocalisation}{وطاقة اللاتمركز} $0.988\beta$.
\end{solution}

\begin{solution}{exo:b2:huckel:11}
في المربع، يجلس الإلكترونان الأعلى طاقة واحدًا في كل مدار من زوج
غير رابط متساوي الطاقة. ومطّ رابطتين متقابلتين يرفع التساوي في الطاقة:
فينخفض أحد المدارين ويرتفع الآخر، ويقترن الإلكترونان كلاهما في الأدنى، مما
يخفض الطاقة. فيستقر الجزيء مستطيلًا برابطتين مزدوجتين
متمركزتين؛ ويبقى \omterm{def:b2:huckel:aromatic}{ضد عطري} وشديد التفاعل.
\end{solution}

\begin{solution}{exo:b2:huckel:12}
المصفوفة (بوحدة $\beta$، مع $\alpha$ أصلًا) \babelsublr{$\begin{pmatrix}0 & 1\\ 1 & 1\end{pmatrix}$}:
$m = 1.618$ و$-0.618$، $E = \alpha + 1.618\beta$ (رابط) و$\alpha - 0.618\beta$. ومن أجل
\omterm{def:b2:diatomic-mos:bonding}{المدار الرابط} $c_X = 1.618c_C$، إذن $c_C^2 = 0.276$، $c_X^2 = 0.724$؛ ومع إلكترونين،
$q_C = 0.55$، $q_X = 1.45$: فكثافة $\pi$ تنجذب نحو X.
\end{solution}

\begin{solution}{pb:b2:huckel:1}
\textbf{1.} $-123.1 - (-4.3) = \qty{-118.8}{kJ/mol}$.
\textbf{2.} $-123.1 - 104.6 = \qty{-227.7}{kJ/mol}$.
\textbf{3.} $-123.1 - 82.9 = \qty{-206.0}{kJ/mol}$.
\textbf{4.} $3 \times (-118.8) = \qty{-356.4}{kJ/mol}$.
\textbf{5.} $356.4 - 206.0 = \qty{150}{kJ/mol}$.
\textbf{6.} يطلق الديين $2 \times 118.8 - 227.7 = \qty{10}{kJ/mol}$ أقل من رابطتين مزدوجتين
معزولتين: فترافق السلسلة المفتوحة يعطي ربحًا صغيرًا، والحلقة المغلقة
للبنزين خمسة عشر ضعفًا منه.
\textbf{7.} مصفوفة $6 \times 6$ فيها 1 بين كل ذرة وجاريها (1--2، 2--3،
\dots، 6--1) و0 في ما عدا ذلك.
\textbf{8.} $m = 2\cos(2\pi k/6)$: $2$، $1$، $1$، $-1$، $-1$، $-2$: $E = \alpha + 2\beta$، $\alpha + \beta$ (مرتين)،
$\alpha - \beta$ (مرتين)، $\alpha - 2\beta$.
\textbf{9.} إلكترونان في $\alpha + 2\beta$، وأربعة في الزوج $\alpha + \beta$.
\textbf{10.} $E_\pi = 6\alpha + 8\beta$.
\textbf{11.} $3(2\alpha + 2\beta) = 6\alpha + 6\beta$.
\textbf{12.} $2\beta$، أي استقرار قدره $2|\beta|$.
\textbf{13.} $2 + 1 + 1 - 1 - 1 - 2 = 0$: مجموع المستويات الستة $6\alpha$.
\textbf{14.} $2|\beta| = \qty{150}{kJ/mol}$: $|\beta| = \qty{75}{kJ/mol}$.
\textbf{15.} $0.472 \times 75 = \qty{35}{kJ/mol}$، مقابل \qty{10}{kJ/mol}: فالنموذج الموفَّق
على البنزين يبالغ في تقدير السلسلة المفتوحة.
\textbf{16.} يمكن أخذ المدارات المشغولة $\psi_1 = (1, 1, 1, 1, 1, 1)/\sqrt6$،
$\psi_2 = (2, 1, -1, -2, -1, 1)/\sqrt{12}$ و$\psi_3 = (0, 1, 1, 0, -1, -1)/2$. ومن أجل الرابطة
1--2: $p_{12} = 2(\frac16 + \frac{2}{12} + 0) = \frac23$؛ وبحكم التناظر لكل رابطة $p = 2/3$.
\textbf{17.} بين مؤشري الرابطتين الطرفية (0.894) والمركزية (0.447) في البوتاديين.
\textbf{18.} الروابط الست متكافئة بحكم التناظر؛ وطولها، \qty{139.7}{pm}، يقع
بين الرابطة المزدوجة للإيثين (\qty{133.9}{pm}) والرابطة الأحادية للإيثان
(\qty{153.6}{pm}).
\textbf{19.} $0.988|\beta| = \qty{74}{kJ/mol}$، نصف طاقة البنزين: فإغلاق الحلقة يضاعف
\omterm{def:b2:huckel:delocalisation}{طاقة اللاتمركز}.
\textbf{20.} $\alpha + 2\beta$، $\alpha + 1.414\beta$ (مرتين)، $\alpha$ (مرتين)، $\alpha - 1.414\beta$ (مرتين)،
$\alpha - 2\beta$.
\textbf{21.} يذهب الإلكترونان الأخيران واحدًا في كل مدار من الزوج غير الرابط: طبقة
مفتوحة، \omterm{def:b2:huckel:aromatic}{ضد عطري}.
\textbf{22.} ينطوي على شكل حوض: فلا تعود مدارات p للروابط المزدوجة المتجاورة
متوازية، وتتناوب الروابط، ويسلك الجزيء سلوك بوليين.
\textbf{23.} أربعة إلكترونات: الزوج نصف الممتلئ نفسه، والحال أسوأ لأن الجزيء لا يستطيع
تجنب الاستواء؛ فيتشوه إلى مستطيل ويكون شديد التفاعل.
\textbf{24.} يكون النظام المترافق الأحادي الحلقة المستوي \omterm{def:b2:huckel:aromatic}{عطريًا} مع $4N + 2$ إلكترون $\pi$،
\omterm{def:b2:huckel:aromatic}{وضد عطري} مع $4N$.
\textbf{25.} $\boldsymbol{|\beta| \approx \qty{75}{kJ/mol}}$.
\end{solution}
